\section{Introduction and main result}\label{sec:intro}

For $\bd=(d_1,\ldots,d_k)$, consider the multihomogeneous polynomial
\begin{equation}\label{eq:target}
F(x,y)=\prod_{i=1}^k x_i^{d_i-1}y_i.
\end{equation}
The pair $(x_i,y_i)$ represents the $i$th symmetric tensor block, of order
$d_i$. Thus $F$ represents a product of $k$ generalized W-states, up to a
nonzero normalization scalar. In this polynomial description, a rank-one
term with respect to the prescribed blocks is a product
$\prod_i(a_i x_i+b_i y_i)^{d_i}$.

\begin{definition}[Partially symmetric rank]
The partially symmetric rank $R_{\bd}(F)$ is the least integer $r$ for which
there are complex coefficients $c_\ell,a_{i\ell},b_{i\ell}$ satisfying
\begin{equation}\label{eq:decomposition}
F(x,y)=\sum_{\ell=1}^r c_\ell
\prod_{i=1}^k(a_{i\ell}x_i+b_{i\ell}y_i)^{d_i}.
\end{equation}
The blocks $(x_i,y_i)$ and their degrees $d_i$ are fixed throughout.
\end{definition}

Canino, Casarotti and Santarsiero~\cite[Theorem~1.1]{CCS2025} give an
explicit decomposition proving the general upper bound
\begin{equation}\label{eq:upper}
R_{\bd}(F)\le M,
\qquad M:=2^{k-1}\left(\sum_{i=1}^k d_i-2k+2\right),
\qquad k\ge2,\quad d_i\ge3.
\end{equation}
Our contribution is the matching lower bound for every decomposition
of the form \eqref{eq:decomposition}, without any restriction on its
support. This is the partially symmetric rank question in problem
TR-29 of the Open Problems in Numerical Linear Algebra catalogue
\cite{TR29}.

\begin{theorem}[Exact partially symmetric rank]\label{thm:main}
Let $k\ge2$ and $d_i\ge3$ for every $i$. Over $\C$,
\begin{equation}\label{eq:main}
R_{\bd}\!\left(\prod_{i=1}^k x_i^{d_i-1}y_i\right)
=2^{k-1}\left(\sum_{i=1}^k d_i-2k+2\right).
\end{equation}
\end{theorem}

The proof of the lower bound is developed in Sections~\ref{sec:coordinates}--\ref{sec:tests}
and completed at the end of Section~\ref{sec:tests}. The upper bound is
\eqref{eq:upper}.

\begin{example}[Three cubic factors]\label{ex:rank20}
For $F=x_1^2y_1x_2^2y_2x_3^2y_3$, the theorem gives
$R_{(3,3,3)}(F)=20$. More generally,
$R_{(3,\ldots,3)}(F)=2^{k-1}(k+2)$.
\end{example}

The proof follows a concrete sequence. An arbitrary $r$-term
decomposition gives $r$ weighted points after a change of coordinates.
Coefficient matching determines their moments. We then construct $M$
test functions and prove that their evaluation vectors are linearly
independent by showing that a weighted Gram matrix is invertible.
Consequently $r\ge M$. The test functions and the matrix blocks are given
explicitly in Section~\ref{sec:tests}; the longer coefficient calculations
are collected in the appendices.

All ranks in \Cref{thm:main} refer to the prescribed partially symmetric
blocks. The argument does not assert the same lower bound for ordinary
tensor rank or border rank. All pairings below are complex bilinear,
using transpose rather than conjugate transpose.

Section~\ref{sec:binary} extends the lower-bound argument to all
products of binary monomials $\prod_i x_i^{a_i}y_i^{b_i}$ with
$a_i\ge b_i\ge1$ and determines their exact partially symmetric rank,
together with a matching explicit decomposition.
